\section{Condensed MPC Optimization Problem}

The previous sections provide the two elements needed for the finite-horizon MPC problem: the condensed prediction model and the lifted constraints. Using \eqref{eq:cap3:stacked-state-prediction}, the predicted state sequence can be eliminated from the cost function, so the optimization problem can be written only in terms of the predicted input sequence \highlight{$\mathbf{\hat{u}}[k]$}.

Starting from

\noindent
\begin{equation}
	J = \norma{x[k]}{P} + \norma{\mathbf{\hat{x}}[k]}{\mathbf{Q}} + \norma{\mathbf{\hat{u}}[k]}{\mathbf{R}},
\end{equation}

\noindent
and replacing

\noindent
\begin{equation}
	\mathbf{\hat{x}}[k] = \mathbf{A}x[k] + \mathbf{B}\mathbf{\hat{u}}[k],
\end{equation}

\noindent
the cost becomes

\noindent
\begin{equation}
	J(\mathbf{\hat{u}}[k]) =
	\mathbf{\hat{u}}[k]^T H \mathbf{\hat{u}}[k]
	+ 2 f[k]^T \mathbf{\hat{u}}[k]
	+ c_0[k],
\end{equation}

\noindent
where

\noindent
\begin{equation}
	H = \mathbf{R} + \mathbf{B}^T \mathbf{Q} \mathbf{B},
	\qquad
	f[k] = \mathbf{B}^T \mathbf{Q} \mathbf{A}x[k],
\end{equation}

\noindent
and

\noindent
\begin{equation}
	c_0[k] = \norma{x[k]}{P} + \norma{\mathbf{A}x[k]}{\mathbf{Q}}.
\end{equation}

The term \highlight{$c_0[k]$} does not depend on the optimization variable and therefore does not affect the optimizer. It is kept here only to show the complete cost expression.

Combining the condensed cost with the lifted constraints, the MPC problem solved at time \highlight{$k$} is

\noindent
\begin{equation}
	\begin{aligned}
		\mathbb{P}(x[k]): \quad
		\min_{\mathbf{\hat{u}}[k] \in \mathbb{R}^{pN}}
		&\quad
		\mathbf{\hat{u}}[k]^T H \mathbf{\hat{u}}[k]
		+ 2 f[k]^T \mathbf{\hat{u}}[k] \\
		\text{subject to}
		&\quad
		\mathbf{S}\mathbf{\hat{u}}[k] \le \mathbf{b}(x[k]).
	\end{aligned}
\end{equation}

\noindent
The matrix \highlight{$H$} defines the quadratic part of the objective, \highlight{$f[k]$} depends on the current state, and the inequality constraints enforce admissible predicted states and inputs over the horizon.

Let \highlight{$\mathbf{\hat{u}}^*[k]$} be the optimal solution of \highlight{$\mathbb{P}(x[k])$}. In receding-horizon control, only the first input of this optimal sequence is applied to the plant:

\noindent
\begin{equation}
	u[k] = \hat{u}^*[k|k].
\end{equation}

\noindent
At the next sampling instant, the state is measured again and the optimization problem is solved again with the updated state.
